\section{The Morse lower bound and the deficit}
\label{sec:counting}

We finish by comparing the count from Proposition~\ref{prop:cap-transfer}
with the ordinary Morse number.  The lower bound uses the number of
generators of the fundamental group, not the augmentation ideal.

\begin{lemma}
\label{lem:morse-lower}
Let $Y$ be a closed connected smooth manifold of even dimension
$n\geq4$, and let $d(\pi_1Y)$ denote the minimum number of generators
of its fundamental group.  Then
\[
 \Morse(Y)\geq\chi(Y)+4d(\pi_1Y).
\]
\end{lemma}

\begin{proof}
For a Morse function on $Y$, let $m_i$ be its number of index-$i$
critical points.  The zero- and one-handles form a connected graph:
handles of index at least two cannot join distinct components.
This graph has $m_0$ vertices and $m_1$ edges, so its fundamental
group is free of rank $m_1-m_0+1$.  Adding the higher handles gives
a surjection from this group onto $\pi_1Y$.  Thus
\[
 d(\pi_1Y)\leq m_1-m_0+1\leq m_1.
\]
Apply the same argument to the negative of the Morse function to
obtain $m_{n-1}\geq d(\pi_1Y)$.  Since $n$ is even, the Euler
identity gives
\[
 \sum_{i=0}^{n}m_i
 =\chi(Y)+2\sum_{i\text{ odd}}m_i
 \geq\chi(Y)+2(m_1+m_{n-1})
 \geq\chi(Y)+4d(\pi_1Y).
\]
Taking the minimum over all Morse functions proves the claim.
\end{proof}

\begin{proof}[Proof of Theorem~\ref{thm:main}]
Choose $b=2$ and $G=A_5^k$ with $k>60^{16b}$ as in
Lemma~\ref{lem:group-gap}.  In particular $d(G)>16b$.
Every finite group is finitely presented.  Gompf's existence theorem
therefore provides a closed connected symplectic four-manifold
$(X,\omega_X)$ with $\pi_1(X)\cong G$ \cite[Theorem~3.1]{Gompf1992}.
The finiteness of $G$ implies $\chi(X)\geq2$, as explained in
Section~\ref{sec:relative-model}.  Lemmas~\ref{lem:augmentation}
and~\ref{lem:row}, followed by Proposition~\ref{prop:relative-morse},
give a Morse function on $X\times\R^8$, equal to $Q$ for $r\geq3$,
with
\[
 N=\chi(X)+4b
\]
critical points.

Proposition~\ref{prop:cap-transfer} produces a smooth one-periodic
Hamiltonian $H$ on the closed connected symplectic manifold
$M=X\times(S^2)^4$.  Its time-one map is nondegenerate at every
fixed point and has exactly
\begin{equation}
 \#\Fix(\phi_H^1)=16N
 =16\chi(X)+64b=\chi(M)+128.
 \label{eq:final-fixed-count}
\end{equation}
Here $\dim_\R M=12$, $\pi_1(M)\cong G$, and
$\chi(M)=16\chi(X)$.  Lemma~\ref{lem:morse-lower} yields
\[
 \Morse(M)\geq\chi(M)+4d(G)
 >\chi(M)+64b=\#\Fix(\phi_H^1).
\]
The cap-transfer proposition also proves that every fixed-point
trajectory of this Hamiltonian is contractible. Hence
\[
 \#\Fix_0(\phi_H^1;H)=\#\Fix(\phi_H^1)<\Morse(M),
\]
which completes the proof.
\end{proof}

\begin{proof}[Proof of Corollary~\ref{cor:deficit}]
Given a real number $D>0$, choose an integer
$m>16b+D/4$ and then a finite integer $k>60^m$.
Lemma~\ref{lem:group-gap} gives $d(A_5^k)>m$.
The bound $b=2$ in Lemma~\ref{lem:augmentation} is independent of
$k$, and all handle operations above take place in the same
dimension twelve with eight added variables.  Applying the
construction for this group and using \eqref{eq:final-fixed-count}
gives
\[
 \Morse(M)-\#\Fix(\phi_H^1)
 \geq4d(A_5^k)-64b>D.
\]
The manifold $X$, and hence $M$, may depend on $D$, as may the
small time and sphere-area parameters used in the construction.
The dimension, quadratic index, and bound $b$ do not depend on $D$.
\end{proof}

In these examples the integral homology is also easy to identify
to the extent needed for the comparison with homological bounds.
The group $G$ is perfect, so $H_1(X;\Z)=0$.
Poincar\'e duality gives $H_3(X;\Z)=H^1(X;\Z)=0$;
the universal coefficient theorem and duality give
$H_2(X;\Z)\cong H^2(X;\Z)\cong\Hom(H_2(X;\Z),\Z)$,
so $H_2(X;\Z)$ is free abelian.  Hence the integral homology of
$M=X\times(S^2)^4$ is torsion-free and concentrated in even degrees.
Its total rank is $\chi(M)$, and the fixed-point count in
\eqref{eq:final-fixed-count} exceeds this number.  Thus the total
fixed-point count exceeds both the rational Betti number and the
rank-and-torsion quantity discussed in the introduction.
The obstruction detected by the
ordinary Morse number is the additional cost of the fundamental group.
